\documentclass[11pt]{amsart}
\usepackage[margin=1in]{geometry}
\usepackage{amsmath,amssymb,amsthm,mathtools}
\usepackage{booktabs,longtable,array,microtype}
\usepackage[hidelinks]{hyperref}
\hypersetup{pdftitle={Square moduli, sparse masks, and sign obstructions for Erdos sums},pdfauthor={Keio}}
\newtheorem{theorem}{Theorem}[section]
\newtheorem{lemma}[theorem]{Lemma}
\newtheorem{proposition}[theorem]{Proposition}
\newtheorem{corollary}[theorem]{Corollary}
\theoremstyle{definition}
\newtheorem{definition}[theorem]{Definition}
\newtheorem{problem}[theorem]{Problem}
\theoremstyle{remark}
\newtheorem{remark}[theorem]{Remark}
\newcommand{\Z}{\mathbb Z}
\newcommand{\Q}{\mathbb Q}
\newcommand{\Zt}{\mathbb Z_{(2)}}
\newcommand{\Ok}{\mathcal O}
\DeclareMathOperator{\rad}{rad}
\DeclareMathOperator{\supp}{supp}
\numberwithin{equation}{section}
\begin{document}
\title[Square moduli and vanishing Erd\H{o}s sums]{Square moduli, sparse masks, and sign obstructions\\for vanishing Erd\H{o}s sums}
\author{Keio}
\date{24 October 2026 --- research draft}
\begin{abstract}
For an odd integer $q>1$, an Erd\H{o}s function modulo $q$ takes values
$\pm1$ away from multiples of $q$ and vanishes on those multiples.
We study the difference between vanishing of its harmonic sum and the
existence of odd vectors in the rational Okada space.
We prove that an odd Okada vector exists if and only if $q$ is a square.
The proof eliminates prime supports in decreasing order on the maximal
$2$-power quotients of unit groups, preserving explicit $2$-adic congruences.
It replaces the finite obstruction computations in the preceding manuscript
by a general argument. More generally, for every prime $\ell$, a rational
Okada vector congruent to $1$ modulo $\ell$ away from zero exists exactly
when $q=m^\ell$ with $m\equiv1\pmod\ell$.

For actual sign vectors, we give exact sparse masks that enlarge the known
counterexamples. Every modulus $3^{2a}5^{2b}7^{2c}$ with $a,b,c\ge1$
carries a vanishing Erd\H{o}s function; this gives a complete classification
on prime support contained in $\{3,5,7\}$. A further mask permits an arbitrary
even power of $11$ to be appended. In contrast, a short analytic argument
proves nonvanishing for every Erd\H{o}s function modulo $195^2$.
We also prove that every odd square satisfies the bounded integer shadow
relaxation, and give an infinite family on which no invariant sign solution
exists. Thus the remaining square-modulus problem is a genuine lifting
problem. The nonvanishing assertion for nonsquare moduli was announced in
2025 by Bharadwaj, Murty, and Pathak; priority for the stronger odd-vector
statement requires comparison with their full manuscript.
\end{abstract}
\maketitle
\tableofcontents

\section{Introduction and relation to the preceding manuscript}
Put $L(1,f)=\sum_{n\ge1}f(n)/n$. For a periodic function, convergence is
equivalent to a zero period sum. An \emph{Erd\H{o}s function modulo $q$}
has $f(0)=0$ and $f(n)\in\{-1,1\}$ for $q\nmid n$. Its zero set forces
$q$ to be its least positive period. Convergence forces $q$ to be odd.
Call $q$ a \emph{counterexample modulus} if such a function has $L(1,f)=0$.
The preceding manuscript \cite{Keio} constructs counterexamples beginning
with $105^2$, and supplies exact certificates at $165^2$, $315^2$,
$495^2$, and $735^2$. This article leaves that manuscript unchanged.

The prime-adic extension gives an exact congruence classification beyond
the binary alphabet: for every prime $\ell$ and modulus $q>1$, a rational
Okada vector taking values in $1+\ell\Z_{(\ell)}$ off zero exists if
and only if $q=m^\ell$ with $m\equiv1\pmod\ell$. This theorem concerns
a prescribed residue class, with no restriction to values $\pm1$.

The main distinction is between three questions: whether an odd rational
Okada vector exists, whether the retained character equations admit the
orbit sums of a sign function, and whether an actual sign function satisfies
the full equations. The first two questions have the same simple answer:
the modulus must be a square, and every odd square works. The third question
has a more restrictive answer. In particular,
\[
 195^2=38025\quad\hbox{is a square with}\quad
 2\varphi(38025)+1=37441<38025,
\]
but it is not a counterexample modulus. This disproves the possible
sufficiency of the elementary totient condition among odd squares.

Our positive constructions use sparse periodic functions which already
have vanishing harmonic sum. Inserting an earlier counterexample into their
zero positions gives independent exponent increases. These are stronger
than the previously known closure under multiplying two counterexample
moduli: the multipliers $25$, $49$, $81$, and $121$ are not themselves
counterexample moduli.

\subsection{Literature and the scope of novelty}
The necessary equations used here are due to Okada \cite{Okada}; their
modern rational-space formulation is given by Murty and Pathak \cite{MP},
and primitive-character formulations are developed by Bharadwaj \cite{Bha}.
The original conjecture is recorded by Livingston \cite{Liv}. For the
algebraic periodic coefficients considered here, a convergent value is
either zero or transcendental, by the established logarithmic-form results
\cite{BBW,MS}.

The announcement of Siddhi Pathak's seminar on 1 October 2025 states that
joint work with A. Bharadwaj and M. R. Murty settles the Erd\H{o}s
nonvanishing question except at perfect-square periods \cite{Seminar}.
The corresponding work is listed as submitted on Pathak's publication
page \cite{Submitted}. We did not obtain its full text. Consequently,
the nonsquare sign-function conclusion is not claimed here as a new
resolution of an open problem. Our proof establishes the stronger
classification of odd $\Zt$-valued Okada vectors; whether that stronger
statement or the same argument occurs in the submitted work remains to
be checked. The square-modulus results below are stated independently
of that unpublished manuscript.

\section{Okada equations and the retained character system}
The Okada equations below hold for any modulus $q>1$; in the binary
discussion we take $q$ odd. Write $U_d=(\Z/d\Z)^\times$ and let $M(q)$ be the
multiplicative monoid generated by the primes dividing $q$. A rational
periodic function with $f(0)=0$ and mean zero belongs to the Okada space
if and only if its harmonic sum vanishes. We use the following standard
forms of its defining equations \cite{Okada,MP,Bha}:
\begin{align}
 &\sum_{m\in M(q)}\frac{f(am)}m=0 &&(a\in U_q),\label{eq:unit}\\
 &\sum_{d\mid q,\ d>1}
 \left(v_p(q/d)-\frac{[p\mid d]}{p-1}\right)S_d=0
 &&(p\mid q),\label{eq:prime}
\end{align}
where $f_d(x)=f((q/d)x)$ and $S_d=\sum_{x\in U_d}f_d(x)$.
The primitive-character form has, for every nonprincipal primitive
Dirichlet character $\chi$ of conductor $m\mid q$, the equation
\begin{equation}\label{eq:character}
 \sum_{m\mid d\mid q}\frac d{\varphi(d)}
 \prod_{p\mid d,\ p\nmid m}\left(1-\frac{\chi(p)}p\right)
 \sum_{x\in U_d}f_d(x)\overline{\chi(x)}=0.
\end{equation}
Together with \eqref{eq:prime} and the mean equation, these are an
equivalent criterion. We make no assertion that every displayed equation
is linearly independent.

Set $R=\Zt$, the rationals with odd denominator. An \emph{odd Okada vector}
means a rational function in this space with $f(0)=0$ and
$f(n)\in1+2R$ for every nonzero residue. Clearing its odd denominators
gives an equivalent odd-integer existence question, without a height bound.
Let $H_d$ be the odd-order subgroup of $U_d$ and put
\[
 Q_d=U_d/H_d=\prod_{p\mid d}Q_p,\qquad
 Q_p\simeq C_{N_p},\quad N_p=2^{v_2(p-1)},\qquad h_d=|H_d|.
\]
For any function define its actual orbit sums
\begin{equation}\label{eq:orbit}
 \Phi_d(y)=\sum_{x\in y}f_d(x),\qquad y\in Q_d.
\end{equation}
If $f$ is odd $R$-valued, each $\Phi_d(y)$ is odd, since $h_d$ is odd.
Characters factoring through $Q_d$ have squarefree conductor. Restricting
\eqref{eq:character} to these characters gives the same equation with
the last sum replaced by $\sum_{y\in Q_d}\Phi_d(y)\overline{\chi(y)}$.
The prime and mean equations also depend only on these sums. We call
this retained system the \emph{$2$-shadow system}.

Conversely, if orbit sums solve the retained system, set
$f_d(x)=\Phi_d(xH_d)/h_d$. Then $f$ is invariant under odd-order units.
Every omitted character coefficient is zero by orthogonality, so the
retained equations imply the full criterion. Since $h_d$ is an odd unit
in $R$, odd orbit sums give an odd Okada vector. This explains exactly
which existence question the shadow decides. For sign functions, the
possible orbit sums have the additional constraints
\begin{equation}\label{eq:bounded-shadow}
 \Phi_d(y)\in\{-h_d,-h_d+2,\ldots,h_d-2,h_d\}.
\end{equation}
Such sums can be lifted to signs within each fibre, but an arbitrary
sign lift need not satisfy the omitted character equations.


% Global theorem; root first-pass audit and independent algebra audit agree.
% Final independent verification records belong in the surrounding paper.
\section{A global obstruction for nonsquare moduli}

Let $R=\Zt=\mathbf Z_{(2)}$. For every odd prime $p$ put
$N_p=2^{v_2(p-1)}$, so that $N_p/(p-1)$ is a unit of $R$.
We use the shadow criterion of the preceding section.

\begin{theorem}\label{thm:global-odd-obstruction}
An odd integer $q>1$ admits an odd rational Okada vector if and only if
$q$ is a perfect square.
\end{theorem}

Both directions are proved below. In particular, this
theorem classifies odd vectors in the rational relaxation; it does not
assert that every odd square admits a vanishing sign function.

\subsection{Compression by prime support}

Write $q=\prod_{p\in P}p^{e_p}$. For every nonempty subset $S\subseteq P$
put $Q_S=\prod_{p\in S}Q_p$, where $Q_p$ is the cyclic group of order
$N_p$. All shadow layers whose prime support is $S$ have this same
domain. Define their sum
\[
 G_S=\sum_{\substack{D\mid q\\\operatorname{supp}(D)=S}}\Phi_D
 \quad\text{on }Q_S.
\]
Since the summands are odd and their number is $\prod_{p\in S}e_p$,
the function $G_S$ has constant parity
\begin{equation}\label{eq:global-parity}
 G_S(x)\equiv\prod_{p\in S}e_p\pmod{2R}.
\end{equation}
In particular, $G_{\{p\}}(x)\equiv e_p\pmod{2R}$.

For a character $\chi$ primitive on $Q_M$, where
$\varnothing\ne M\subseteq P$, extend $\chi$ to each $Q_S$ with
$S\supseteq M$ by projection and write
$\widehat G_S(\chi)=\sum_{x\in Q_S}G_S(x)\overline{\chi(x)}$.
Dividing the character equation by
$\prod_{p\in M}p/(p-1)$ gives
\begin{equation}\label{eq:global-A}
 \sum_{S\supseteq M}
 \prod_{p\in S\setminus M}\frac{p-\chi(p)}{p-1}
 \widehat G_S(\chi)=0.
\end{equation}

For every $p\in P$ define
\[
 T_p(G)=\sum_{S\ni p}\sum_{x\in Q_S}G_S(x).
\]
The mean equation converts the original $p$-equation to
\[
 \sum_{\substack{D\mid q\\p\mid D}}
 \bigl((p-1)v_p(D)+1\bigr)\sum_{x\in Q_D}\Phi_D(x)=0.
\]
Every layer sum belongs to $2R$, since $Q_D$ has even cardinality
and every entry of $\Phi_D$ is odd. Consequently
\begin{equation}\label{eq:global-T}
 T_p(G)=-(p-1)\sum_{\substack{D\mid q\\p\mid D}}
 v_p(D)\sum_{x\in Q_D}\Phi_D(x)\in 2N_pR.
\end{equation}

\subsection{An integral elimination step}

\begin{lemma}[Decomposition with controlled parity]
\label{lem:global-decomp}
Let $|S|\ge2$ and let $F:Q_S\to R$ have constant parity. If its
Fourier coefficients vanish for every character nontrivial on each
factor $Q_p$, $p\in S$, then there are functions
$h_p:Q_{S\setminus\{p\}}\to R$ of constant parity such that
\[
 F(x)=\sum_{p\in S}h_p(x_{S\setminus\{p\}}).
\]
\end{lemma}

\begin{proof}
Fix the identity element as base point in each factor. The anchored
inclusion--exclusion decomposition is
\[
 F(x)=\sum_{U\subseteq S}F_U(x_U),\qquad
 F_U(x_U)=\sum_{V\subseteq U}(-1)^{|U|-|V|}
 F(x_V,0_{S\setminus V}).
\]
The vanishing hypothesis says that $F$ is a sum, over the rationals,
of functions missing at least one coordinate. It therefore implies
$F_S=0$. Equivalently, this follows by applying the full mixed
difference to the Fourier expansion.
For $U\ne\varnothing$, all $2^{|U|}$ terms in $F_U$ have the same
parity, so $F_U$ is even-valued. The term $F_\varnothing=F(0)$ is
constant. Assign each proper subset $U\subsetneq S$ to some
$p\in S\setminus U$ and inflate $F_U$ to
$Q_{S\setminus\{p\}}$. Summing the terms assigned to each $p$
gives the required $h_p$. Each is even-valued except possibly the one
receiving the constant term, and that one also has constant parity.
\end{proof}

\begin{lemma}[Elimination circuit]\label{lem:global-circuit}
Suppose $S\subseteq P$, $|S|\ge2$, $p\in S$, and
$h:Q_{S\setminus\{p\}}\to R$ has constant parity. Write
$U=S\setminus\{p\}$. The changes
\begin{equation}\label{eq:global-circuit}
 \delta G_S=-\operatorname{inf}_{U}^{S}h,
 \qquad
 \delta G_U=\frac{N_p}{p-1}(p-T_p^{\mathrm{tr}})h
\end{equation}
preserve every character equation \eqref{eq:global-A}, every congruence
\eqref{eq:global-T}, and the parity of $G_U$. Here
$(T_p^{\mathrm{tr}}h)(y)=h(y+\overline p)$ is translation by the image
of $p$ in $Q_U$.
\end{lemma}

\begin{proof}
The coefficient $N_p/(p-1)$ belongs to $R$. Since $p$ is odd and
$h$ has constant parity, $(p-T_p^{\mathrm{tr}})h$ is even-valued.
Thus the lower-layer change is even and preserves its parity.

For a character whose primitive support contains $p$, the change on
$Q_S$ has zero Fourier coefficient, because it is constant in the
$Q_p$ coordinate; the lower layer does not contribute. For primitive
support $M\subseteq U$, inflation multiplies its Fourier coefficient
by $N_p$, while
\[
 \widehat{T_p^{\mathrm{tr}}h}(\chi)=\chi(p)\widehat h(\chi).
\]
The two changes in \eqref{eq:global-A}, after removing the common
factor indexed by $U\setminus M$, are therefore
\[
 -\frac{p-\chi(p)}{p-1}N_p\widehat h(\chi)
 +\frac{N_p}{p-1}(p-\chi(p))\widehat h(\chi)=0.
\]
All other primitive supports are unaffected.

Put $H=\sum_{x\in Q_U}h(x)$. Translation preserves sums, so the
total changes in the two layers are $-N_pH$ and $+N_pH$.
Consequently the total mean is unchanged, and $T_s(G)$ is unchanged
for every $s\ne p$. The change in $T_p(G)$ is $-N_pH$.
The set $U$ is nonempty, so $|Q_U|$ is even; since $h$ has constant
parity, $H\in2R$. Thus $-N_pH\in2N_pR$, as required.
\end{proof}

\subsection{Proof of the global obstruction}

\begin{proof}[Proof of Theorem~\ref{thm:global-odd-obstruction}, necessity]
Suppose that an odd Okada vector exists. Form the compressed functions
$G_S$ above, which satisfy \eqref{eq:global-A} and
\eqref{eq:global-T} and have the constant parities
\eqref{eq:global-parity}.

Process the subsets $S\subseteq P$ of size at least two in decreasing
order of their size. When $S$ is reached, all strict supersets have
already been eliminated. Hence \eqref{eq:global-A} for primitive
support $S$ says precisely that $G_S$ has no fully primitive Fourier
component. Every earlier alteration of this layer was even-valued,
so it still has constant parity. Lemma~\ref{lem:global-decomp}
expresses it as a sum of functions missing one coordinate. Apply
Lemma~\ref{lem:global-circuit} to each summand. These operations
replace $G_S$ by zero, alter only smaller layers, preserve every
character equation and every congruence \eqref{eq:global-T}, and
preserve the parity of each smaller layer. The process terminates.

Only the singleton functions $G_{\{p\}}$ remain. Their nonprincipal
Fourier coefficients vanish by \eqref{eq:global-A}, so each is a
constant $a_p\in R$. Its parity is still $e_p$ modulo $2R$.
But \eqref{eq:global-T} now gives
\[
 N_pa_p=T_p(G)\in2N_pR,
\]
so $a_p$ is even, and therefore $e_p$ is even. This holds for every
$p\mid q$. Thus $q$ is a perfect square.
\end{proof}

\begin{proof}[Proof of Theorem~\ref{thm:global-odd-obstruction}, sufficiency]
Suppose every $e_p$ is even. For each nonempty support $S\subseteq P$
choose one prime $p(S)\in S$ and a nontrivial quadratic character
$\eta_{p(S)}$ on $Q_{p(S)}$. On each layer $D$ with support $S$, set
\[
 \Phi_D(y)=(-1)^{v_{p(S)}(D)}\eta_{p(S)}(y_{p(S)}).
\]
Every value is $1$ or $-1$, and each layer sum is zero. Hence the mean
and all prime equations vanish. For each fixed support $S$, summing
over the even number $e_{p(S)}$ of possible positive exponents of
$p(S)$ gives $G_S=0$. Thus all character equations vanish as well.
Shadow descent yields an odd rational Okada vector. Its original
layer values are $\Phi_D/|H_D|$; their denominators are odd, but they
need not have absolute value one.
\end{proof}

\begin{corollary}
Every convergent Erd\H{o}s sign sum at an odd nonsquare modulus is
nonzero, and hence transcendental. Any vanishing Erd\H{o}s sign
function has an odd square modulus.
\end{corollary}

The maximal $2$-power quotient has a precise role in the argument:
$N_p/(p-1)$ must be a unit of $\Zt$ in the elimination circuit.
Replacing $Q_p$ by a smaller $2$-quotient can introduce an even
denominator and invalidate the preservation of parity.

% Priority caveat for the surrounding manuscript: the 2025 announced
% Bharadwaj--Murty--Pathak nonsquare theorem must be compared in full.
% Their announced sign-function theorem does not alone state the stronger
% odd-rational-vector classification, but its proof may imply it.

\section{Sparse zero-sum masks and inflation of counterexamples}
\label{sec:sparse-inflation}

The usual composition of two counterexamples multiplies their moduli.
The construction below also permits multiplication by certain prime squares
which themselves admit no counterexample.  Its auxiliary functions have
several zero residues, and therefore are not themselves Erd\H{o}s functions.

\begin{definition}
Let $1<m\mid Q$. A \emph{zero-harmonic mask of type $(Q,m)$} is a
$Q$-periodic function $g\colon\mathbb Z\to\{-1,0,1\}$ such that
\[
 g(n)=0\quad\Longleftrightarrow\quad m\mid n,
 \qquad \sum_{n=1}^{Q}g(n)=0,
 \qquad L(1,g)=0.
\]
\end{definition}

\begin{theorem}[Mask insertion]
\label{thm:mask-insertion}
Suppose that $f$ is an Erd\H{o}s counterexample of modulus $q$, and that
$g$ is a zero-harmonic mask of type $(Q,m)$ with $Q\mid mq$. Then
\[
 H(n)=
 \begin{cases}
 g(n),&m\nmid n,\\
 f(n/m),&m\mid n
 \end{cases}
\]
is an Erd\H{o}s counterexample of modulus $mq$. In particular, a mask of
type $(am,m)$ permits multiplication by $m$ of every counterexample modulus
divisible by $a$.
\end{theorem}

\begin{proof}
The condition $Q\mid mq$ makes $H$ periodic modulo $mq$. Its zero set is
exactly $mq\mathbb Z$, and all other values are signs. The identity
\[
 H(n)=g(n)+\boldsymbol{1}_{m\mid n}f(n/m)
\]
shows that its sum over a period is zero. Both series on the right converge,
so the same identity gives
\[
 L(1,H)=L(1,g)+\frac1m L(1,f)=0.
\]
Moreover, its least positive period is $mq$, as is forced by its zero set.
No coprimality hypothesis is needed.
\end{proof}

\begin{proposition}[Three exact masks]
\label{prop:three-masks}
There exist zero-harmonic masks of each of the three types in
Table~\ref{tab:sparse-masks}.
\end{proposition}

\begin{table}[ht]
\centering
\begin{tabular}{r r r l r}
\hline
$Q$ & $m$ & $Q/m$ & Character orders used & Unit equations \\
\hline
$11025$ & $49$ & $225$ & $2,2,2$ at $3,5,7$ & $5040$ \\
$11025$ & $25$ & $441$ & $2,4,2$ at $3,5,7$ & $5040$ \\
$99225$ & $81$ & $1225$ & $2,2,2$ at $3,5,7$ & $45360$ \\
\hline
\end{tabular}
\caption{Certified masks. In each case all the indicated unit equations,
the three prime equations, and the period sum vanish exactly.}
\label{tab:sparse-masks}
\end{table}

\begin{proof}
The divisor-layer sign tables are the supplementary certificates
\texttt{mask\_11025\_hole\_49.json},
\texttt{mask\_11025\_hole\_25.json}, and
\texttt{mask\_99225\_hole\_81.json}.
For each divisor $d>1$ of $Q$, write
$n=(Q/d)x$ with $(x,d)=1$. A table row gives $g(n)$ as a function of the
discrete logarithm classes of $x$ at the primes dividing $d$, listed in
increasing order. The coordinates have the orders shown in
Table~\ref{tab:sparse-masks}. The primitive roots are $2,2,3$ modulo
$3,5,7$, respectively. A missing row is identically zero.

The independent standard-library program
\texttt{verify\_sparse\_masks.py} expands these tables to every residue,
checks the advertised zero sets and signs, and verifies Okada's original
equations using exact rational arithmetic. For each $p\mid Q$, it computes
the resolvent
\[
 T_p h(a)=\sum_{j\ge0}\frac{h(p^j a)}{p^j}
\]
by solving $T_p h(a)=h(a)+T_p h(pa)/p$ on every cycle and incoming tree of
the map $a\mapsto pa\pmod Q$. Every local recurrence is checked again.
The commuting product of these resolvents gives the smooth sum for every
unit residue. All $\varphi(Q)$ such values vanish. The three integer-scaled
prime-valuation equations and the period sum also vanish. Consequently
Okada's criterion gives $L(1,g)=0$ in each case. This verification does not
use the projected shadow equations or the integer-programming solver used
to discover the tables.
\end{proof}

\begin{corollary}[Three universal inflation rules]
\label{cor:inflation-rules}
If $q$ is a counterexample modulus, then so are
\[
 \begin{array}{ll}
 49q,&\text{if }225\mid q,\\
 25q,&\text{if }441\mid q,\\
 81q,&\text{if }1225\mid q.
 \end{array}
\]
\end{corollary}

\begin{theorem}[All square moduli on the radical $105$]
\label{thm:all-357-squares}
For every $a,b,c\ge1$, the modulus
\[
 q=3^{2a}5^{2b}7^{2c}
\]
admits an Erd\H{o}s counterexample.
\end{theorem}

\begin{proof}
The previously certified seeds $105^2=3^25^27^2$ and
$315^2=3^45^27^2$ give the cases $a=1$ and $a=2$, with $b=c=1$.
Every constructed modulus remains divisible by $225$, $441$, and $1225$.
The three rules of Corollary~\ref{cor:inflation-rules} therefore allow us
to increase $c$ by $1$, $b$ by $1$, and $a$ by $2$, independently.
The two seeds cover the two parities of $a$.
\end{proof}

\begin{corollary}[Classification of odd $\{3,5,7\}$-smooth moduli]
\label{cor:357-classification}
Let $q>1$ be odd and have no prime factors outside $\{3,5,7\}$. Then $q$
is a counterexample modulus if and only if it is a square divisible by
$105^2$.
\end{corollary}

\begin{proof}
Theorem~\ref{thm:all-357-squares} proves sufficiency. The non-square
obstruction proved in this paper excludes every non-square $q$ with this
prime support. If a square $q$ omits at least one of the three primes, then
\[
 \frac{\varphi(q)}q\ge
 \left(1-\frac13\right)\left(1-\frac15\right)
 =\frac8{15}>\frac12.
\]
Proposition~\ref{prop:mask-support}, applied to an Erd\H{o}s function,
would require $2\varphi(q)\le q-1$, excluding this case.
\end{proof}

The inflation statements also apply to counterexamples with additional
prime factors. For example, the seed $165^2$ gives all moduli
$165^2\cdot49^k$, $k\ge0$.

\subsection{A necessary support bound}

The next estimate explains why the smaller multiplier $9$ cannot be obtained
from a mask of period $11025$.

\begin{proposition}[Support bound]
\label{prop:mask-support}
Let $h$ be a rational $Q$-periodic function with $h(0)=0$, zero period sum,
$L(1,h)=0$, and $|h(a)|=1$ for every unit $a$ modulo $Q$. Then
\[
 2\varphi(Q)\le\sum_{n\bmod Q}|h(n)|.
\]
In particular, a zero-harmonic mask of type $(Q,m)$ can exist only if
\[
 2\varphi(Q)\le Q-Q/m.
\]
\end{proposition}

\begin{proof}
Write $M(Q)$ for the monoid generated by the primes dividing $Q$.
For each unit $a$, the unit Okada equation and the triangle inequality give
\[
 1=|h(a)|\le
 \sum_{\substack{r\in M(Q)\\r>1}}\frac{|h(ar)|}{r}.
\]
Every residue $n\pmod Q$ has total weight exactly $1$ in the double sum
\[
 \sum_{a\in(\mathbb Z/Q\mathbb Z)^\times}
 \sum_{r\in M(Q)}\frac{|h(ar)|}{r}.
\]
Indeed, if $d=(n,Q)$, solutions of $ar\equiv n\pmod Q$ with $a$ a unit
exist exactly when $(r,Q)=d$, and their number is
$\varphi(Q)/\varphi(Q/d)$. To evaluate the resulting weight prime by
prime, let $p^e\Vert Q$ and $v=\min(v_p(n),e)$. If $v<e$, the local
fibre size is $p^v$ and the sole exponent in $r$ is $v$, giving weight
$p^vp^{-v}=1$. If $v=e$, the local weight is
\[
 \varphi(p^e)\sum_{j\ge e}p^{-j}
 =(p-1)p^{e-1}\frac{p^{1-e}}{p-1}=1.
\]
All terms are nonnegative, so the rearrangement is justified. Summing the
unit inequalities and removing the $r=1$ contribution gives
$\varphi(Q)\le\sum_n|h(n)|-\varphi(Q)$, as required.
\end{proof}

For $Q=11025$, a mask with $m=9$ would have support size
$11025-1225=9800$, while $2\varphi(11025)=10080$. It is therefore
impossible, without any symmetry restriction. More generally, on radical
$105$, the support inequality forces
\[
 m\ge\frac{1}{1-2\varphi(Q)/Q}=\frac{35}{3}.
\]
Among positive even powers of $3$, the first permitted multiplier is $81$.
Thus the use of an $81$-mask and two initial $3$-exponents reflects a genuine
support obstruction.

\subsection{An explicit 49-mask}
\label{subsec:49-mask-table}

For completeness, the shortest certificate is displayed in
Table~\ref{tab:49-mask}. For $n=(11025/d)x$ with $(x,d)=1$, each row lists
the signs in lexicographic order of the quadratic-character coordinates at
the primes dividing $d$. Coordinate $0$ means Legendre symbol $+1$ and
coordinate $1$ means $-1$. The last coordinate varies fastest. Every omitted
divisor layer is zero. The listed layers are exactly those with $7\mid d$,
so the zero set is exactly the multiples of $49$.

\begin{table}[ht]
\centering
\begin{tabular}{r l@{\qquad}r l@{\qquad}r l}
\hline
$d$ & Signs & $d$ & Signs & $d$ & Signs \\
\hline
$7$ & \texttt{+-} & $147$ & \texttt{-+-+} & $735$ & \texttt{+--++--+}\\
$21$ & \texttt{-+-+} & $175$ & \texttt{+--+} & $1225$ & \texttt{+--+}\\
$35$ & \texttt{-++-} & $245$ & \texttt{+--+} & $1575$ & \texttt{+--++--+}\\
$49$ & \texttt{+-} & $315$ & \texttt{++-+---+} & $2205$ & \texttt{+--++--+}\\
$63$ & \texttt{--++} & $441$ & \texttt{--++} & $3675$ & \texttt{+--++--+}\\
$105$ & \texttt{++----++} & $525$ & \texttt{+-+++---} & $11025$ & \texttt{-++--++-}\\
\hline
\end{tabular}
\caption{A zero-harmonic mask of type $(11025,49)$.}
\label{tab:49-mask}
\end{table}

\subsection{Further masks on the radical 165}

Two additional compact certificates give zero-harmonic masks of types
\[
 (Q,m)=(27225,121),\qquad (Q,m)=(245025,81).
\]
Their primitive-root coordinates have orders $2,4,2$ at $3,5,11$,
respectively, with primitive root $2$ at all three primes. The same independent
full-residue verifier checks $13200$ and $118800$ unit equations,
respectively, as well as all prime equations, the exact zero sets, and the
zero period sums. The certificates are
\texttt{mask\_27225\_hole\_121.json} and
\texttt{mask\_245025\_hole\_81.json}.

\begin{corollary}
Every counterexample modulus $q$ has the following further inflation rules:
\[
 \begin{aligned}
 225\mid q&\ \Longrightarrow\ 121q\text{ is a counterexample modulus},\\
 3025\mid q&\ \Longrightarrow\ 81q\text{ is a counterexample modulus}.
 \end{aligned}
\]
In particular, every modulus $3^{2a}5^2 11^{2c}$, $a,c\ge1$, admits a
counterexample. More generally, the same holds for
\[
 3^{2a}5^{2b}11^{2c}\qquad(a,c\ge b\ge1).
\]
\end{corollary}

\begin{proof}
Apply the mask insertion theorem. The previously certified seeds
$165^2=3^25^211^2$ and $495^2=3^45^211^2$, together with the multipliers
$81$ and $121$, give all exponents with $b=1$. A product of $b$ such
counterexample moduli gives any prescribed $a,c\ge b$, by writing each of
$a$ and $c$ as a sum of $b$ positive integers and using the existing
composition rule for ordinary counterexamples.
\end{proof}

Combining the $121$ inflation rule with
Theorem~\ref{thm:all-357-squares} also gives every modulus
\[
 3^{2a}5^{2b}7^{2c}11^{2d},
 \qquad a,b,c\ge1,\quad d\ge0.
\]
Thus the fourth prime can have any nonnegative even exponent, independently
of the three positive even exponents in the complete radical-$105$ family.
No claim is made here about the still-uncovered exponents with prime support
$\{3,5,11\}$, nor about an arbitrary additional prime.

\section{The modulus \texorpdfstring{$195^2$}{195 squared}}\label{sec:q195}

The distinction between an odd Okada vector and a sign vector is already
visible at $q=195^2$.  We prove the conjectured nonvanishing at this modulus
for arbitrary Erd\H{o}s functions, without an invariance hypothesis.
The proof combines a bound on the reciprocal mass in Okada's unit equations
with three character equations of conductor $13$ and the prime-$3$ equation.

\begin{lemma}[Weighted forcing]\label{lem:weighted-forcing}
Let $q>1$ be odd, let $M(q)$ be the multiplicative monoid generated by its
prime divisors, and put
\[
 R_q=\prod_{p\mid q}\frac{p}{p-1},\qquad
 E_q=\frac{R_q(1-1/q)-2}{2}.
\]
Suppose that an Erd\H{o}s function $f$ modulo $q$ has $L(1,f)=0$.
For every $a\in U_q$ the reciprocal mass of the agreeing terms is
\[
 \sum_{\substack{m\in M(q),\ m>1,\ q\nmid m\\ f(am)=f(a)}}\frac1m=E_q.
\]
Consequently, if $m\in M(q)$, $m>1$, $q\nmid m$ and $1/m>E_q$, then
$f(am)=-f(a)$ for every $a\in U_q$.
\end{lemma}

\begin{proof}
Okada's unit equation is
$\sum_{m\in M(q)}f(am)/m=0$.
The full reciprocal mass is $R_q$, and the terms divisible by $q$ have mass
$R_q/q$ and contribute zero.
After removing $m=1$, the remaining nonzero terms have total mass
$R_q(1-1/q)-1$.  Multiplying the equation by $f(a)$ shows that the mass of
disagreeing terms exceeds that of agreeing terms by $1$.
Solving these two equations gives the asserted value $E_q$.
An individual agreeing term cannot have mass greater than their total mass.
\end{proof}

\begin{theorem}\label{thm:q195-exclusion}
Every convergent Erd\H{o}s sum of modulus $q=195^2=38025$ is nonzero.
\end{theorem}

\begin{proof}
Assume to the contrary that $L(1,f)=0$, and write $q=3^25^213^2$.
Here $R_q=65/32$ and
\[
 E_q=\frac1{64}-\frac1{37440}<\frac1{64}.
\]
Lemma~\ref{lem:weighted-forcing} therefore gives
\begin{equation}\label{eq:195-forced-multipliers}
 f(ma)=-f(a)\qquad(a\in U_q),\qquad
 m\in\{3,5,9,13,15,25,27,39,45\}.
\end{equation}

\smallskip\noindent
\emph{The unit restriction.}
The equation for $m=9$ implies that the restriction of $f$ to $U_q$ depends
only on the image in $U_{q/9}$: if $a\equiv b\pmod{q/9}$, then
$f(a)=-f(9a)=-f(9b)=f(b)$.
Under $U_q\cong U_9\times U_{25}\times U_{169}$ this removes the $U_9$
coordinate.  The equations for $m=25$ and $m=13$ similarly remove the
$U_{25}$ coordinate and the kernel of $U_{169}\to U_{13}$.
Hence $f(a)=G(a\bmod13)$ for $a\in U_q$.

Given $a\in U_q$, choose $b\in U_q$ with
$b\equiv3a\pmod{q/9}$; this is possible by the Chinese remainder theorem,
since $q/9=4225$ is coprime to $3$.
Then $9b\equiv27a\pmod q$, and
\eqref{eq:195-forced-multipliers} gives $f(b)=f(a)$.
It follows that $G(3t)=G(t)$.
The element $3$ has order $3$ modulo $13$, and $2$ is a primitive root with
$\langle3\rangle=\{2^0,2^4,2^8\}$.
There are thus signs $F_j\in\{\pm1\}$, indexed by $j\in\mathbb Z/4\mathbb Z$,
such that
\begin{equation}\label{eq:195-unit-pattern}
 f(a)=F_{\operatorname{ind}_2(a\bmod13)\bmod4}\qquad(a\in U_q).
\end{equation}

Put $D=\{1,3,5,9,13,15,25,39,45\}$.
For $g\in D$, set $d=q/g$.  Every $x\in U_d$ lifts to a unit modulo $q$,
and $13\mid d$.  Equations \eqref{eq:195-forced-multipliers} and
\eqref{eq:195-unit-pattern} consequently determine the entire layer:
\begin{equation}\label{eq:195-forced-layers}
 f(gx)=\sigma_gF_j,\qquad
 \sigma_g=\begin{cases}1,&g=1,\\-1,&g>1,\end{cases}
 \qquad j=\operatorname{ind}_2(x\bmod13)\bmod4.
\end{equation}
The counts used below are recorded in Table~\ref{tab:195-forced}.
Each $j$-fibre in a listed layer has $\varphi(d)/4$ elements.
These layers contain $36552$ nonzero residues, leaving only $1472$ values
of $f$ unrestricted.

\begin{table}[ht]
\centering
\caption{Forced layers for $q=38025$.  The last column is
$c_d=2(2-v_3(d))-[3\mid d]$, the coefficient in twice the prime-$3$ equation.}
\label{tab:195-forced}
\begin{tabular}{rrrrr}
\toprule
$g$ & $d=q/g$ & $\varphi(d)/4$ & $\sigma_g$ & $c_d$\\
\midrule
1  & 38025 & 4680 &  1 & -1\\
3  & 12675 & 1560 & -1 &  1\\
5  &  7605 &  936 & -1 & -1\\
9  &  4225 &  780 & -1 &  4\\
13 &  2925 &  360 & -1 & -1\\
15 &  2535 &  312 & -1 &  1\\
25 &  1521 &  234 & -1 & -1\\
39 &   975 &  120 & -1 &  1\\
45 &   845 &  156 & -1 &  4\\
\bottomrule
\end{tabular}
\end{table}

\smallskip\noindent
\emph{The character equations.}
For every $d\mid q$ with $13\mid d$, define the actual fibre sums
\[
 \Phi_d(j)=
 \sum_{\substack{x\in U_d\\\operatorname{ind}_2(x\bmod13)\equiv j\pmod4}}
 f((q/d)x),\qquad j\in\mathbb Z/4\mathbb Z,
\]
and put
\[
 A_j=\sum_{\substack{d\mid q\\13\mid d,\ 5\nmid d}}\Phi_d(j),\qquad
 B_j=\sum_{\substack{d\mid q\\65\mid d}}\Phi_d(j).
\]
These definitions impose no invariance assumption on $f$.
For each of the three nonprincipal characters of
$U_{13}/\langle3\rangle\cong C_4$, the character equation in Okada's criterion
asserts that the corresponding Fourier coefficient of
\begin{equation}\label{eq:195-H}
 H_j=8A_j+10B_j-2B_{j+1}
\end{equation}
vanishes.  Thus $H_j$ is constant.
Indeed, multiply the layered character equation by
$(\varphi(13)/13)(3-1)(5-1)=96/13$.
Since $\chi(3)=1$, the contribution of the prime $3$ is the factor $2$
whether or not $3\mid d$; the contribution of $5$ is $4$ if $5\nmid d$ and
$5-\chi(5)$ otherwise.
Finally $\operatorname{ind}_2(5)\equiv1\pmod4$, so multiplication of a Fourier
coefficient by $\chi(5)$ is translation $j\mapsto j+1$.
This gives \eqref{eq:195-H}.

By \eqref{eq:195-forced-layers}, the forced part of $A_j$ is $-234F_j$,
and the forced part of $B_j$ is
\[
 (4680-1560-936-780-360-312-120-156)F_j=456F_j.
\]
The total cardinality of a $j$-fibre across the $A$-layers is
\[
 \frac{(\varphi(13)+\varphi(169))(1+\varphi(3)+\varphi(9))}{4}
 =\frac{168\cdot9}{4}=378,
\]
so at most $144$ entries are unrestricted.
Across the $B$-layers the corresponding cardinality is
\[
 \frac{168\cdot9(\varphi(5)+\varphi(25))}{4}=9072.
\]
The forced $B$-fibres contain $4680+4224=8904$ entries, leaving $168$
unrestricted entries.
We may therefore write
\[
 A_j=-234F_j+a_j,\quad |a_j|\le144;\qquad
 B_j=456F_j+b_j,\quad |b_j|\le168.
\]
Equation~\eqref{eq:195-H} becomes
\begin{equation}\label{eq:195-H-bound}
 H_j=2688F_j-912F_{j+1}+e_j,
 \qquad |e_j|\le8\cdot144+12\cdot168=3168.
\end{equation}
If the cyclic sign pattern $F$ is nonconstant, it has both a transition
$+1\to-1$ and a transition $-1\to+1$.
At these two indices the forced expression in
\eqref{eq:195-H-bound} is respectively $3600$ and $-3600$.
Constancy of $H$ would require their difference $7200$ to be canceled by
an error difference of absolute value at most $6336$.
This is impossible.  Consequently $F$ is constant.

\smallskip\noindent
\emph{The prime equation.}
Write $S_d=\sum_{x\in U_d}f((q/d)x)$.
Twice Okada's prime-$3$ equation is
\[
 \sum_{\substack{d\mid q\\d>1}}c_dS_d=0,
 \qquad c_d=2(2-v_3(d))-[3\mid d]\in\{4,1,-1\}.
\]
If $S=\sum_{j=0}^3F_j$, the contribution of the forced layers is
\[
 \begin{split}
 &(-4680-1560+936-3120+360-312+234-120-624)S\\
 &\hspace{35mm}=-8886S.
 \end{split}
\]
The remaining $1472$ signs have coefficients of absolute value at most $4$,
so their total contribution has absolute value at most $5888$.
But a constant $F$ has $S=\pm4$, giving $|8886S|=35544>5888$.
This final contradiction proves the theorem.
\end{proof}

\begin{remark}
The argument uses real size bounds in addition to the $2$-primary character
projection.  It does not require parity of fibre sums, an optimization
solver, or enumeration of sign patterns.
In particular, the existence of bounded odd shadow solutions at a square
modulus does not prevent the weighted unit equations from forcing additional
layer values that make the resulting real bounds inconsistent.
\end{remark}

\section{Why bounded orbit sums do not decide the square case}
\label{sec:bounded}

The global theorem excludes all nonsquares. On every remaining odd square,
even the exact integer and size conditions on individual orbit sums are
feasible. This explains why the additional forcing argument at $195^2$
is essential.

\begin{theorem}[Saturation of the bounded shadow]
\label{thm:bounded-squares}
For every odd square $q>1$, the $2$-shadow equations admit orbit sums
satisfying \eqref{eq:bounded-shadow}. In fact every orbit sum can be
chosen to have absolute value one. Consequently, some sign function
has zero retained shadow, whether or not $q$ is a counterexample modulus.
\end{theorem}
\begin{proof}
Use the direct square construction in the sufficiency proof of
Theorem~\ref{thm:global-odd-obstruction}. Every coordinate $\Phi_d(y)$
there is $1$ or $-1$, and all retained equations hold. The fibre size
\[
 h_d=\prod_{p^a\parallel d}p^{a-1}
              \frac{p-1}{2^{v_2(p-1)}}
\]
is a positive odd integer, so either choice $\Phi_d(y)=\pm1$ satisfies
\eqref{eq:bounded-shadow}. On a fibre with sum $\Phi$ and size $h$,
choosing $(h+\Phi)/2$ positive signs and $(h-\Phi)/2$ negative signs
produces a sign lift. No assertion about the omitted character equations
of that lift is needed.
\end{proof}

\begin{remark}
For fixed feasible orbit sums, the exact number of sign lifts is
\[
 \prod_{d\mid q,\ d>1}\prod_{y\in Q_d}
 \binom{h_d}{(h_d+\Phi_d(y))/2}.
\]
These lifts need not have zero harmonic sum. For example, at $q=9$ the
vector indexed by $0,1,\ldots,8$,
\[
 (0,1,1,1,-1,1,-1,-1,-1),
\]
has zero retained shadow. The primitive character modulo $9$ with
$\chi(2)=e^{2\pi i/6}$, however, gives
$\sum_{a\in U_9}f(a)\overline{\chi(a)}=4$, violating the full
conductor-$9$ equation. Its order is $6$, a direction discarded by the
$2$-shadow. In particular, integer feasibility of the bounded shadow
must never be reported as a sign counterexample.
\end{remark}


% Standalone theorem/proof fragment for the follow-up paper.
% Original counterexample manuscript is not modified.
\section{A height obstruction to odd-order invariance}

The existence of an odd vector in the Okada kernel does not imply the
existence of a sign vector in its invariant part.  The following theorem
gives an infinite family on which this distinction is unavoidable.
Here the height of an integer-valued function is its maximum absolute value.

\begin{theorem}[Large-prime symmetry and height obstruction]
\label{thm:large-prime-symmetry}
Let $N\geq1$ be odd, let $r\nmid N$ be an odd prime, let $e\geq1$, and
put $q=Nr^e$ and $m=2^{v_2(r-1)}$.  Suppose that every prime divisor of
$N$ is an $m$-th power residue modulo $r$.  Under the Chinese remainder
identification
\[
  U_q=U_N\times U_{r^e},
\]
let $K=\{1\}\times H_{r^e}$, where $H_{r^e}$ is the subgroup of odd order
of $U_{r^e}$.

Suppose that $f$ is an integer-valued, $q$-periodic function such that
$f(0)=0$, every $f(n)$ for $n\not\equiv0\pmod q$ is odd, $f$ is invariant
under multiplication by $K$, and $L(1,f)$ converges and vanishes.
Then
\[
  N\,\lVert f\rVert_\infty\geq r.
\]
Consequently, if $r>N$, there is no $K$-invariant Erd\H{o}s function of
modulus $Nr^e$ with vanishing sum.  In particular there is no such
function invariant under the full odd-order subgroup of $U_q$.
\end{theorem}

\begin{proof}
Set $B=\lVert f\rVert_\infty$ and suppose, for a contradiction, that
$NB<r$.  Let $Q=U_r/H_r$; this cyclic group has order $m$.  Reduction
modulo $r$ identifies the quotient $U_{r^j}/H_{r^j}$ with $Q$, for each
$1\leq j\leq e$.

Every residue $n\pmod q$ with $v_r(n)=e-j$ has a unique representation
\[
 n=r^{e-j}u,\qquad u\pmod {Nr^j},\quad r\nmid u.
\]
The $K$-orbits of these residues are indexed by
$(a,y)\in(\mathbb Z/N\mathbb Z)\times Q$, where $a=u\pmod N$ and $y$
is the image of $u\pmod {r^j}$ in $Q$.  Thus $K$-invariance defines
odd integers $b_j(a,y)$ by
\[
 b_j(a,y)=f(r^{e-j}u),\qquad
 \Gamma_j(y)=\sum_{a\pmod N}b_j(a,y).
\]
Every fibre $(a,y)$ contains
\[
 h_j=\frac{r^{j-1}(r-1)}m
\]
residues.  Since $N$ is odd, every $\Gamma_j(y)$ is odd; also
$|\Gamma_j(y)|\leq NB<r$.

Let $\chi\ne1$ be a character of $Q$, viewed as a primitive Dirichlet
character of conductor $r$.  The residue hypothesis gives
$\chi(p)=1$ for every prime $p\mid N$.  We compute its layered Okada
coefficient.  Write a divisor containing $r$ as $d=d_0r^j$, with
$d_0\mid N$.  Its weight is
\[
 \frac{d}{\varphi(d)}
 \prod_{p\mid d_0}\left(1-\frac{\chi(p)}p\right)
 =\frac r{r-1}.
\]
If $n=(q/d)x=r^{e-j}(N/d_0)x$, then $u=(N/d_0)x$ and
$\chi(u)=\chi(x)$, again by the residue hypothesis.  Summing the layers
over all $d_0\mid N$ therefore sums precisely the residues in the
$j$-th band above.  The character equation becomes
\begin{align*}
 0
 &=\frac r{r-1}\sum_{j=1}^e h_j
      \sum_{y\in Q}\Gamma_j(y)\overline{\chi(y)}\\
 &=\frac1m\sum_{y\in Q}
       \left(\sum_{j=1}^e r^j\Gamma_j(y)\right)
       \overline{\chi(y)}.
\end{align*}
The necessity of these equations is the layered Okada criterion.
Fourier orthogonality on $Q$ shows that
\[
 P(y)=\sum_{j=1}^e r^{j-1}\Gamma_j(y)
\]
is independent of $y$.

For $y,z\in Q$, the integers
$a_j=(\Gamma_j(y)-\Gamma_j(z))/2$ satisfy $|a_j|\leq NB<r$ and
$\sum_{j=1}^e r^{j-1}a_j=0$.  Reduction modulo $r$ gives $a_1=0$;
division by $r$ and repetition gives every $a_j=0$.  Thus
\[
 \Gamma_j(y)=c_j\quad\hbox{for all }y\in Q,
\]
where each $c_j$ is odd and $|c_j|\leq NB<r$.

The remaining nonzero residues are $n=r^e a$, $1\leq a<N$.  Put
\[
 S_0=\sum_{a=1}^{N-1}f(r^e a),\qquad
 |S_0|\leq B(N-1)<r-1.
\]
The strict final inequality follows because $NB$ and $r$ are integers,
so $NB\leq r-1$ and $B\geq1$.
Convergence of a periodic Dirichlet series at $1$ implies mean zero,
which here reads
\[
 0=S_0+\sum_{j=1}^e mh_jc_j
   =S_0+(r-1)\sum_{j=1}^e r^{j-1}c_j.
\]
Consequently $r-1$ divides $S_0$, so $S_0=0$ and
$\sum_{j=1}^e r^{j-1}c_j=0$.  Reduction modulo $r$ gives $c_1=0$,
because $|c_1|<r$, contradicting its oddness.
\end{proof}

\begin{remark}
The proof needs only mean zero and the $m-1$ conductor-$r$ character
equations.  It does not use any prime-logarithm equation, including
the prime-$r$ equation.  No converse is claimed: these few equations are necessary,
not a sufficient criterion for membership in the full Okada kernel.
The theorem concerns odd integer vectors with a height restriction;
it does not prohibit odd $2$-integral rational vectors of arbitrary
height or denominator.
\end{remark}

\begin{remark}[Sharpness when $N=1$]
The height bound is attained at $q=r^2$.  Let $\chi$ be the quadratic
character modulo $r$ and define
\[
 f(n)=-\chi(n)\quad(r\nmid n),\qquad
 f(ra)=r\chi(a)\quad(r\nmid a),\qquad f(0)=0.
\]
This is an odd integer vector of height $r$, invariant under
$H_{r^2}$.  Both layer means are zero.  The only nonzero character
coefficients occur at $\chi$, where their weighted sum is
$r\cdot r+r^2\cdot(-1)=0$.  Hence the layered Okada criterion gives
$L(1,f)=0$ and equality in the bound $\lVert f\rVert_\infty\geq r$.
\end{remark}

\begin{corollary}[Infinitely many squares with no invariant sign solution]
Let $r>11025$ be prime and $r\equiv-1\pmod {420}$.  Set
\[
 q=11025r^2=(105r)^2.
\]
There is no Erd\H{o}s function of modulus $q$ with vanishing sum that is
invariant under the odd-order subgroup of the local $r$-component.
Nevertheless $q$ has an odd Okada vector, by the all-odd-squares theorem,
and
\[
 \frac{\varphi(q)}q=\frac{16}{35}\left(1-\frac1r\right)<\frac12.
\]
There are infinitely many such primes $r$.
\end{corollary}

\begin{proof}
The congruence gives $r\equiv3\pmod4$, so $m=2$.  Also
$r\equiv11\pmod {12}$ gives $(3/r)=1$, while quadratic reciprocity gives
\[
 (5/r)=(r/5)=(-1/5)=1,\qquad
 (7/r)=-(r/7)=-(-1/7)=1.
\]
The theorem applies with $N=11025$ and $e=2$.  Since
$\gcd(419,420)=1$, Dirichlet's theorem supplies infinitely many primes
in this progression.
\end{proof}

The first prime $r>11025$ in this progression is $13859$, giving
$q=2117592488025$.  The result does not rule out a non-invariant
Erd\H{o}s counterexample at these moduli.  Rather, any counterexample
there would have to have a nonzero Fourier component on which the
odd-order local subgroup acts nontrivially.  It therefore marks a
provable limitation of restricting the sign search to the $2$-shadow.

% Full prime-adic classification, independently audited by three agents.
% Uses the support-elimination lemmas in global_obstruction_candidate.tex.
% Novelty relative to the submitted Bharadwaj--Murty--Pathak work is unverified.
\section{A complete prime-adic congruence classification}

Let $\ell$ be a prime and put $R_\ell=\mathbb Z_{(\ell)}$.
Congruences in $R_\ell$ refer to its residue field.

\begin{theorem}\label{thm:ell-classification}
For an integer $q>1$, the following conditions are equivalent:
\begin{enumerate}
\item There is a rational-valued function $f$ of period $q$ such that
\[
 f(0)=0,\quad f(n)\in1+\ell R_\ell\ (q\nmid n),\quad
 \sum_{a=1}^qf(a)=0,\quad L(1,f)=0.
\]
\item $q\equiv1\pmod\ell$ and $\ell\mid v_p(q)$ for every prime $p\mid q$.
\item $q=m^\ell$ for an integer $m\equiv1\pmod\ell$.
\end{enumerate}
\end{theorem}

For $\ell=2$ this is precisely the odd-square classification.
For larger $\ell$ it concerns a rational congruence class of coefficients,
not the binary sign alphabet. The equivalence of the last two conditions
uses $m^\ell\equiv m\pmod\ell$.

\begin{proof}[Necessity]
The mean equation reduced modulo $\ell$ gives $q-1\equiv0\pmod\ell$;
in particular $\ell\nmid q$. Write $q=\prod_pp^{e_p}$ and partition its
prime divisors into
\[
 P_+=\{p:\ell\mid p-1\},\qquad P_0=\{p:\ell\nmid p-1\}.
\]
Set $N_p=\ell^{v_\ell(p-1)}$ and let $Q_p$ have order $N_p$;
thus $Q_p$ is trivial for $p\in P_0$. These are the maximal
$\ell$-quotients of the local unit groups, since $p\ne\ell$.
For $D\mid q$, the fibre size is
$m_D=\varphi(D)/\prod_{p\mid D}N_p$.

Average under the subgroup of units of order prime to $\ell$ and form
the fibre sums $\Phi_D$. They satisfy the Okada equations, and
$\Phi_D(y)\equiv m_D\pmod\ell$. Their totals are
$S_D=\sum_y\Phi_D(y)\equiv\varphi(D)\pmod\ell$.
As before, the mean equation converts the prime equation to
\[
 B_p^*:=\sum_{D:p\mid D}((p-1)v_p(D)+1)S_D=0.                 \tag{\thesection.1}
\]
The elementary divisor-sum identity
\[
 \sum_{D:p\mid D}((p-1)v_p(D)+1)\varphi(D)=e_p(p-1)q
\]
shows immediately that $\ell\mid e_p$ for every $p\in P_0$.

It remains to treat $P_+$. Write $G_S$ for the sum of the functions
$\Phi_D$ with support $S$, on $Q_S=\prod_{p\in S}Q_p$.
With $h_p=(p-1)/N_p$ for $p\in P_+$,
its constant residue is
\[
 G_S\equiv
 \prod_{p\in S\cap P_+}e_ph_p
 \prod_{b\in S\cap P_0}(b^{e_b}-1)\pmod\ell.                 \tag{\thesection.2}
\]
Indeed, the sum of the positive-exponent fibre sizes at a prime in
$P_+$ is $e_ph_p$ modulo $\ell$, whereas at $b\in P_0$ it is
$\sum_{j=1}^{e_b}\varphi(b^j)=b^{e_b}-1$.
For $p\in P_+$, every $S_D$ with $p\mid D$ belongs to $\ell R_\ell$.
Equation (1) therefore implies
\[
 T_p(G):=\sum_{S\ni p}\sum_yG_S(y)\in\ell N_pR_\ell.         \tag{\thesection.3}
\]

Fold the trivial factors into the nontrivial ones. For nonempty
$S\subseteq P_+$ define
\[
 \widetilde G_S=
 \sum_{U\subseteq P_0}
 \left(\prod_{b\in U}\frac{b-\tau_b}{b-1}\right)G_{S\cup U},
 \qquad (\tau_bF)(y)=F(y+\overline b).                        \tag{\thesection.4}
\]
All operators are $R_\ell$-integral, commute, preserve total sums,
and act as the identity on constant residues. Consequently the
normalized character equations regroup exactly into the same system
on $P_+$, and $T_p(\widetilde G)=T_p(G)$ for $p\in P_+$.
Writing $r=\prod_{b\in P_0}b^{e_b}$, (2) yields
\[
 \widetilde G_S\equiv r\prod_{p\in S}e_ph_p\pmod\ell.         \tag{\thesection.5}
\]

Apply the support-elimination argument of
Lemmas~\ref{lem:global-decomp} and~\ref{lem:global-circuit}, replacing
parity by constant residue modulo $\ell$. Here are the necessary
checks for this replacement. An anchored difference annihilates a
constant residue, so the decomposition into functions missing a
coordinate is $R_\ell$-integral and each summand has constant residue.
The transfer operator is still
\[
 \frac{N_p}{p-1}(p-\tau_p).
\]
Its coefficient is a unit of $R_\ell$, and its value on a
constant-residue function is zero modulo $\ell$, because
$p\equiv1\pmod\ell$. The remaining coordinate set is nonempty and
has cardinality divisible by $\ell$, so the change in $T_p$ lies in
$\ell N_pR_\ell$. Character cancellation is unchanged.
This descent uses the character equations, (3), and the constant
residues; a separate mean equation for the folded functions is not
needed.

After eliminating all supports of size at least two, each singleton
is a constant $a_p$, with
$a_p\equiv re_ph_p\pmod\ell$ by (5). Equation (3) gives
$N_pa_p\in\ell N_pR_\ell$, hence $a_p\equiv0\pmod\ell$.
Both $r$ and $h_p$ are units, so $\ell\mid e_p$. If $P_+$ is empty,
this step is unnecessary. Necessity follows for every prime divisor.
\end{proof}

\begin{proof}[Sufficiency]
Assume condition 2. First let $p\in P_+$ and $\ell\mid e$.
Choose a period-$p$ function $u$ with $u(0)=0$, $u(1)=2-p$, and
$u(a)=1$ at all other nonzero residues. It has zero mean and all
nonzero-residue values are $1$ modulo $\ell$. Set
\[
 c_t=1\ (0\leq t<e-1),\qquad
 c_{e-1}=-\sum_{t=0}^{e-2}p^{e-1-t}.
\]
Every $c_t\equiv1\pmod\ell$ and $\sum_tc_tp^{-t}=0$.
The period-$p^e$ function with value $c_tu(n/p^t)$ when $v_p(n)=t<e$,
and zero when $p^e\mid n$, has zero mean and
\[
 L(1,g)=\left(\sum_{t=0}^{e-1}\frac{c_t}{p^t}\right)L(1,u)=0.
\]

Now put $r=\prod_{b\in P_0}b^{e_b}$. If $r=1$, omit this factor.
Otherwise $r\equiv1\pmod\ell$, since every prime in $P_+$ is
$1$ modulo $\ell$. Every $\varphi(D)$ for $D\mid r$ is a unit of
$R_\ell$, and all its quotient groups are trivial. We construct a
function constant on each gcd layer. Start with layer totals
$\Phi_D=\varphi(D)$ for $D>1$. Their mean is $r-1$ and their
prime equations are $B_b^*=e_b(b-1)r$, all divisible by $\ell$.

Choose a prime $b_0\mid r$ and subtract $r-1$ from $\Phi_{b_0}$.
The mean is now zero and the new prime totals are
\[
 \beta_b=e_b(b-1)r-\mathbf1_{b=b_0}b(r-1)\in\ell R_\ell.
\]
For each prime $b\mid r$, set $w_b=-\beta_b/(b-1)\in\ell R_\ell$;
subtract $w_b$ from $\Phi_b$ and add it to $\Phi_{b^2}$.
The layers $b^2$ exist because $e_b\geq\ell\geq2$.
This preserves the mean, changes $B_b^*$ by $(b-1)w_b=-\beta_b$,
and leaves every other prime equation unchanged. All resulting
$\Phi_D$ remain congruent to $\varphi(D)$ modulo $\ell$.

Assign the constant value $\Phi_D/\varphi(D)$ to the layer with
denominator $D$. All nonprincipal character coefficients vanish,
because the function is constant on each unit layer; the mean and
prime equations have just been made zero. Thus the full vanishing
criterion gives $L(1,g)=0$, while every nonzero-residue value is
$1$ modulo $\ell$.

Finally compose the constructed factors: for periods $A,B$, use
$F(n)=u(n)$ if $A\nmid n$ and $F(n)=v(n/A)$ otherwise.
The period is $AB$, the mean remains zero, the congruence class is
preserved, and
$L(1,F)=L(1,u)+A^{-1}L(1,v)$. This constructs the required function
for $q$.
\end{proof}

\paragraph{Two explicit ternary-congruence examples.}
For $\ell=3$, $q=7^3$, the first construction uses
$u(1)=-5$, $u(2)=\cdots=u(6)=1$ and $(c_0,c_1,c_2)=(1,1,-56)$.
Its Dirichlet-series multiplier is $1+1/7-56/49=0$.
The second construction applies to $q=64=4^3$, although its prime
divisor $2$ is not $1$ modulo $3$. Set
\[
 f(0)=0,\qquad f(32)=196,\qquad f(16)=f(48)=-128,
\]
and set $f(a)=1$ at every other nonzero residue modulo $64$.
All nonzero-residue values are $1$ modulo $3$ and the period sum is zero. Every
unit equation has value
$1+1/2+1/4+1/8-128/16+196/32=0$; its prime-2 equation is also zero.
Neither example is a binary sign counterexample.

\section{Verification, consequences, and remaining questions}

The nonsquare obstruction and the exclusion at $195^2$ are analytic
proofs; neither depends on optimization output or a finite testing range.
Their algebra was checked separately, including the character translations
and the preservation of the congruences used in the elimination.
The sign constructions have a different status: they are explicit finite
certificates verified by exact rational arithmetic. The supplement includes
compact divisor-layer tables, full value arrays, and standard-library
programs that check the original unit and prime equations independently
of the discovery model. No small floating-point $L$-value is used to
certify equality to zero.

The original seeds required for the construction are included as data,
with exact full-equation checks. A second implementation checks the sparse
masks using functional graphs of multiplication on all residues. For
$195^2$, independent programs check the elementary counts and also rebuild
the original equation space before checking its separating certificates.
For the global theorem, exact coefficient-level tests of the elimination
circuits supplement the proof; they are not a substitute for its induction.
These are internal independent implementations and mathematical reviews,
not external peer review.

The earlier finite computations on $255$ exponent cores supported on
$\{3,5,7,11\}$ are retained in the supplement as regression data.
They are no longer needed to infer a theorem about nonsquares. Likewise,
the constructive exponent-reduction argument is supplied as an additional
note: each positive exponent may be reduced to $1$ if odd and $2$ if even
without changing odd-vector existence. The global classification makes
that equivalence immediate, but the explicit transformations may be useful
for manipulating certificates.

Three questions from the preceding manuscript now have more precise
answers. Its general odd-vector conjecture is proved; the question at
$195^2$ has a negative answer; and sparse masks construct infinite families
beyond product closure. The suggestion that every odd square satisfying
$2\varphi(q)+1\le q$ might work is false, by $q=195^2$.
Determining all counterexample square moduli remains open in this work.
In particular, the large-prime symmetry obstruction does not exclude
non-invariant sign functions at the moduli it describes. A search which
restricts to invariant signs cannot settle those cases.

The same fixed sign alphabet at higher integer arguments does not supply
a further open nonvanishing problem: for $k\ge2$,
\[
 |L(k,f)|\ge 1-\sum_{\substack{n\ge2\\q\nmid n}}n^{-k}
 =2-(1-q^{-k})\zeta(k)>0.
\]
This is the elementary first-term estimate, not a new result about the
Chowla--Milnor conjecture. No unrelated named open conjecture is claimed
to be resolved here.

\begin{thebibliography}{99}
\bibitem{BBW} A. Baker, B. J. Birch, and E. A. Wirsing,
On a problem of Chowla, \emph{J. Number Theory} \textbf{5} (1973), 224--236.
\bibitem{Bha} A. Bharadwaj,
On primitivity and vanishing of Dirichlet series,
arXiv:2204.03674 (2022), \url{https://arxiv.org/abs/2204.03674}.
\bibitem{Keio} Keio,
Erd\H{o}s's conjecture on $\sum f(n)/n$ is false:
a $2$-adic shadow of Okada's criterion and the least counterexample
$q=105^2$, preceding unpublished manuscript and supplementary data,
24 October 2026.
\bibitem{Liv} A. E. Livingston,
The series $\sum_1^\infty f(n)/n$ for periodic $f$,
\emph{Canad. Math. Bull.} \textbf{8} (1965), 413--432.
\bibitem{MP} M. R. Murty and S. Pathak,
The Okada space and vanishing of $L(1,f)$,
\emph{Funct. Approx. Comment. Math.} \textbf{66} (2022), 35--57,
\url{https://www.cmi.ac.in/~siddhi/Okada_space_RM_SP_Final.pdf}.
\bibitem{MS} M. R. Murty and N. Saradha,
Transcendental values of the digamma function,
\emph{J. Number Theory} \textbf{125} (2007), 298--318.
\bibitem{Okada} T. Okada,
On a certain infinite series for a periodic arithmetical function,
\emph{Acta Arith.} \textbf{40} (1982), 143--153.
\bibitem{Seminar} S. Pathak,
Seminar announcement, 1 October 2025,
Number Theory Seminar 2025--26, Charles University;
joint work with A. Bharadwaj and M. R. Murty,
\url{https://sites.google.com/view/shman/number-theory-seminar-2025-26},
accessed 24 October 2026.
\bibitem{Submitted} A. Bharadwaj, M. R. Murty, and S. Pathak,
Dirichlet series with periodic coefficients and a conjecture of Erd\H{o}s,
listed as submitted on S. Pathak's publication page, item 21,
\url{https://www.cmi.ac.in/~siddhi/publications.html},
accessed 24 October 2026. Full text not consulted.
\end{thebibliography}
\end{document}
